\SourcePage{196}{199}
\section{The uncertainty relation for energy}

Consider a system made up of two weakly interacting parts. Suppose that at a
certain instant these parts are known to have definite energies, denoted by
$E$ and $\varepsilon$. After an interval $\Delta t$, the energies are measured
again, giving values $E'$ and $\varepsilon'$ which generally differ from $E$
and $\varepsilon$. We can readily determine the order of magnitude of the
most probable value of $E'+\varepsilon'-E-\varepsilon$ found by the
measurement.

According to (42.3), with $\omega=0$, the probability that a time-independent
perturbation transfers the system during time $t$ from energy $E$ to energy
$E'$ is proportional to
\[
  \frac{\sin^2\left(\dfrac{E'-E}{2\hbar}t\right)}{(E'-E)^2}.
\]
\SourcePage{197}{200}
It follows that the most probable difference $E'-E$ is of order $\hbar/t$.

Applying this result to the present case, with the interaction between the
two parts as the perturbation, gives
\begin{equation}
  |E+\varepsilon-E'-\varepsilon'|\Delta t\sim\hbar.
  \tag{44.1}
\end{equation}
Thus the shorter the interval $\Delta t$, the larger the observed energy
change. Its order of magnitude, $\hbar/\Delta t$, is importantly independent
of the strength of the perturbation.

The energy change specified by (44.1) will be observed however weak the
interaction between the two parts may be. This is a purely quantum result
with a profound physical meaning. It shows that in quantum mechanics two
measurements can test energy conservation only to an accuracy of order
$\hbar/\Delta t$, where $\Delta t$ is the time between them.

Relation (44.1) is often called the uncertainty relation for energy. It must,
however, be emphasized that its meaning is essentially different from that
of the coordinate--momentum relation $\Delta p\Delta x\sim\hbar$. In the
latter, $\Delta p$ and $\Delta x$ are the uncertainties in momentum and
coordinate at one and the same instant; the relation says that these two
quantities cannot have strictly definite values simultaneously.

The energies $E$ and $\varepsilon$, by contrast, can be measured at each
given instant with arbitrary precision. The quantity
$(E+\varepsilon)-(E'+\varepsilon')$ in (44.1) is the difference between two
exactly measured energy values at two different times, not an uncertainty in
the energy at a particular instant.

If $E$ is regarded as the energy of a system and $\varepsilon$ as the energy
of a ``measuring apparatus,'' their interaction energy can be taken into
account only to an accuracy $\hbar/\Delta t$. Denote the errors in the
measurements of the corresponding quantities by
$\Delta E,\Delta\varepsilon,\ldots$. In the favorable case where
$\varepsilon$ and $\varepsilon'$ are known exactly
($\Delta\varepsilon=\Delta\varepsilon'=0$),
\begin{equation}
  \Delta(E-E')\sim\frac{\hbar}{\Delta t}.
  \tag{44.2}
\end{equation}

This relation has important consequences for momentum measurement. Measuring
the momen%
\SourcePageJoin{198}{201}
tum of a particle (take an electron for definiteness) involves a
collision with another, ``measuring,'' particle whose momenta before and
after the collision may be regarded as known
exactly.\footnote{For the present argument, how the energy of the
``measuring'' particle becomes known is immaterial.} Applying momentum
conservation to this collision gives three equations, the components of one
vector equation, for six unknowns: the components of the electron momentum
before and after the collision. More equations might be sought by arranging
a sequence of collisions with measuring particles and applying momentum
conservation to each. But the number of unknowns also grows, since it includes
the electron momenta between collisions; for any number of collisions there
are plainly three more unknowns than equations. Momentum measurement must
therefore use energy conservation at each collision in addition to momentum
conservation. As we have seen, the energy law can be applied only to an
accuracy of order $\hbar/\Delta t$, with $\Delta t$ the interval between the
beginning and end of the process under consideration.

For simplicity, consider an idealized thought experiment in which the
``measuring particle'' is a perfectly reflecting plane mirror. Only the
momentum component normal to the mirror then matters. Momentum and energy
conservation give, for determining the particle momentum $P$,
\begin{equation}
  p'+P'-p-P=0,
  \tag{44.3}
\end{equation}
\begin{equation}
  |\varepsilon'+E'-\varepsilon-E|\sim\frac{\hbar}{\Delta t}
  \tag{44.4}
\end{equation}
($P,E$ are the particle's momentum and energy; $p,\varepsilon$ are those of
the mirror; unprimed and primed quantities refer respectively to times before
and after the collision). The quantities $p,p',\varepsilon,\varepsilon'$ of
the ``measuring particle'' may be treated as known exactly, so their errors
vanish. The equations above then give
\[
  \Delta P=\Delta P',
  \qquad
  |\Delta E'-\Delta E|\sim\frac{\hbar}{\Delta t}.
\]
But
\[
  \Delta E=\frac{\partial E}{\partial P}\Delta P=v\Delta P,
\]
\SourcePage{199}{202}
and similarly
\[
  \Delta E'=v'\Delta P'=v'\Delta P.
\]
Hence
\begin{equation}
  |(v_x'-v_x)\Delta P_x|\sim\frac{\hbar}{\Delta t}.
  \tag{44.5}
\end{equation}
The suffix $x$ has been attached to the velocities and momentum to stress
that the relation applies separately to each component.

This is the required relation. It shows that measuring an electron's momentum
to a specified accuracy $\Delta P$ is unavoidably accompanied by a change in
its velocity, and therefore in the momentum itself.

The shorter the measurement process, the greater this change. The velocity
change can be made arbitrarily small only as $\Delta t\to\infty$, but a
momentum measurement extending over a long time can in general have meaning
only for a free particle. Here the nonrepeatability of momentum measurements
at short intervals, and the two-sided nature of measurement in quantum
mechanics, are especially clear: one must distinguish the measured value of
a quantity from the value produced by the measuring process.\footnote{Relation
(44.5), as well as the clarification of the physical meaning of the energy
uncertainty relation, is due to N.~Bohr (1928).}

The argument at the beginning of this section, based on perturbation theory,
can also be approached through the decay of a system under some perturbation.
Let $E_0$ be an energy level calculated while entirely neglecting the
possibility of decay. Denote by $\tau$ the \emph{lifetime} of this state, the
reciprocal of its decay probability per unit time. The same method gives
\begin{equation}
  |E_0-E-\varepsilon|\sim\hbar/\tau,
  \tag{44.6}
\end{equation}
where $E$ and $\varepsilon$ are the energies of the two parts into which the
system has decayed. Their sum $E+\varepsilon$ provides a measure of the
system's energy before decay. The relation therefore shows that the energy of
a decaying system in a \emph{quasistationary} state can be defined only to an
accuracy of order $\hbar/\tau$. This quantity is customarily called the level
\emph{width} $\Gamma$. Thus
\begin{equation}
  \Gamma\sim\hbar/\tau.
  \tag{44.7}
\end{equation}
